

\section{Densidad espectral de potencia de la señal PAM}

Consideramos la onda PAM (tren de pulsos modulados en amplitud)
\[
x(t) = \displaystyle\sum_k a_k p(t-kT_s),
\]
en la cual hacemos las siguientes suposiciones:
\begin{itemize}
    \item El pulso es $p(t)$ arbitrario, de energía finita, con transformada de Fourier $\hat{P}(f)$.
    \item Los símbolos $\{a_k\}$ forman un proceso aleatorio de tiempo discreto, estacionario y ergódico, de densidad
    espectral de potencia $S_a(\varphi)$. Definiremos estos conceptos en breve.
   \end{itemize}

Para estudiar a la se\~nal $x(t)$ mediante herramientas de procesos estacionarios (de tiempo continuo), ser\'a adem\'as
necesario ``des-sincronizar'' el análisis introduciendo un retardo aleatorio, al igual que se hizo en la Secci\'on \ref{sec.estoc}.


\subsection*{Procesos aleatorios de tiempo discreto}

Cubrimos primero los elementos fundamentales de procesos, que no hemos tocado antes pero que son matemáticamente
análogos al caso de tiempo continuo. Una familia de variables aleatorias $\{a_k\}_{k\in \mathbb{Z}}$ es un proceso estacionario en sentido amplio si:
\begin{itemize}
\item $E[a_k]=\mu$ independiente de $k$;
\item $E[a_k a_{k-l}] = R_a(l)$, independiente de $k$.
\end{itemize}
La ergodicidad implica que la autocorrelaci\'on estoc\'astica $R_a(l)$  coincide con probabilidad 1
con la autocorrelación determinística discreta
\begin{align*}
r_{a}(l) & = \lim_{K\to\infty} \frac{1}{2K+1} \sum_{k=-K}^{K}a_k a_{k-l}.
\end{align*}
La anterior autocorrelación fue definida en el curso de Análisis de Señal,
donde también se introdujo su Tranformada de Fourier de Tiempo
Discreto (TFTD), la densidad espectral de potencia discreta
\begin{align}\label{eq.safi}
S_a(\varphi) = \sum_{l=-\infty}^{+\infty} R_a(l) \e^{-j 2\pi\varphi l}.
\end{align}
Esta funci\'on, peri\'odica de per\'iodo 1, depende de la variable de frecuencia normalizada $\varphi$.


\subsection*{Teorema}
Sea $\{a_k\}$ un proceso estacionario de s\'imbolos, con densidad espectral de potencia discreta $
S_a(\varphi)$. Sea  $\theta$ es una V.A. uniforme en
$[0,T_s]$ independiente del proceso de símbolos. Entonces $y(t)$ definido por
\[
 y(t)=x(t+\theta)=\displaystyle\sum_k
a_k p(t+\theta - kT_s)\]
es un proceso estacionario de tiempo continuo, con densidad espectral de potencia

\begin{equation}
   \label{eq.dig}
    {\color{Red}\boxed{\normalcolor S_y(f) = f_s \left|\hat{P}(f)\right|^2
    S_a\left(\frac{f}{f_s}\right)}}%\tag{\#}
\end{equation}

\begin{proof} Calculamos la autocorrelación del proceso $y(t)$. Para ello, para un $\theta$ fijo calculamos el valor esperado
condicional $E\left[y(t) y(t+\tau)|\theta\right]$ (sobre el
proceso $\{a_k\}$), y despues promediamos en $\theta$ uniforme.
\begin{align*}
    E\left[y(t)y(t+\tau)|\theta\right]  & = E\left[\sum_k a_k p(t+\theta-kT_s) \sum_m a_m p(t+\tau+\theta-mT_s)\right]\\
                        & = \sum_{k,m} E[a_k a_m] p(t+\theta - kT_s) p(t+\tau+\theta-mT_s)\\
        \underset{\scriptsize(l=m-k)}{} & =  \sum_{k,l} R_a(l) p(t+\theta-kT_s) p(t+\tau+\theta-lT_s-kT_s)\\
                        & = \sum_l R_a(l) \sum_k p(t+\theta-kT_s)
                        p(t+\tau-lT_s+\theta-kT_s).
\end{align*}

Ahora promedio en $\theta$:
\begin{align}\label{eq.rytau}
    E\left[y(t)y(t+\tau)\right]         & = \frac{1}{T_s}\int_0^{T_s} E\left[y(t)y(t+\tau)|\theta\right]\d\theta  \nonumber \\
                        & = f_s \sum_l R_a(l) \sum_k\int_0^{T_s} p(t+\theta-kT_s) p(t+\theta-kT_s + \tau-lT_s)\d\theta \nonumber\\
\underset{\scriptsize(\eta=\theta-kT_s+t)}{}    & = f_s  \sum_l R_a(l) \sum_k \int_{t-kT_s}^{t-kT_s+T_s} p(\eta) p(\eta + \tau - lT_s)\d\eta\nonumber\\
                        & = f_s \sum_l R_a(l) \int_{-\infty}^{+\infty} p(\eta) p(\eta + \tau-lT_s)\d\eta\nonumber\\
                        & = f_s \sum_l R_a(l) r_p(\tau-lT_s).
\end{align}

En el \'ultimo paso introdujimos
\[
r_p(\tau) := \left\langle p(t),p(t+\tau)\right\rangle = \int_{-\infty}^{+\infty} p(t) p(t + \tau)\d t,
\]
autocorrelación del pulso como señal de {\em energía} finita. Notar que el resultado es independiente de $t$.

Un argumento similar (más sencillo) prueba que $E[y(t)]$ es  independiente de $t$. Por lo tanto $y(t)$ es un proceso continuo
estacionario.

Recordamos del curso de An\'alisis de Se\~nal que la transformada de Fourier de $r_p(\tau)$ es la densidad espectral de energ\'ia
$|\hat{P}(f)|^2$. Transformando por Fourier en (\ref{eq.rytau}) obtenemos
\begin{align*}
    S_y(f)  & = f_s \sum_l R_a(l) \e^{-j2\pi(lT_s)f} \overbrace{\mathcal{F}\left[r_p(\tau)\right]}^{\left|\hat{P}(f)\right|^2}\\
        & = f_s \left|\hat{P}(f)\right|^2\sum_l R_a(l) \e^{-j2\pi l
        (\frac{f}{f_s})}.
\end{align*}
Recurriendo a (\ref{eq.safi}) interpretamos la \'ultima suma como
$S_a\left(\frac{f}{f_s}\right).$

\end{proof}


\subsection*{Caso particular m\'as usado}

En la Secci\'on \ref{sec.densdig} consideramos un caso particular del proceso de s\'imbolos: $\{a_k\}$ independientes, id\'enticamente distribuidos, con $E[a_k]=\mu$, $E[a_k^2]=\mu^2+\sigma^2$.
Veamos ahora como particularizar la f\'ormula (\ref{eq.dig}) a ese caso.

Utilizando la independencia, calculamos la autocorrelaci\'on del proceso de s\'imbolos:
\begin{align*}
R_a(l) &= E[a_k a_{k+l}]= \begin{cases}
                                        \mu^2+\sigma^2 & l = 0\\
                                        \mu^2 & l\neq 0
                                       \end{cases} \\
                                       &= \mu^2+\sigma^2\delta(l).
\end{align*}

Su transformada de Fourier de tiempo discreto (funci\'on peri\'odica de $\varphi\in\R$) es
\[
  S_a(\varphi) = \sigma^2 + \mu^2 \sum_{n=-\infty}^{+\infty}\delta(\varphi - n).
\]

\[\Rightarrow  S_a\left(\frac{f}{f_s}\right) = \sigma^2+\mu^2\sum_{n=-\infty}^{+\infty} \delta\left(\frac{f-nf_s}{f_s}\right) = \sigma^2+\mu^2 f_s \sum_{n=-\infty}^{+\infty} \delta (f-nf_s).\]

Sustiyendo en (\ref{eq.dig}) llegamos a la f\'ormula utilizada en la Secci\'on \ref{sec.densdig}.
\begin{equation}
    \label{eq:dig3}
    {\color{Red}\boxed{\normalcolor S_x(f) = \sigma^2 f_s \left|\hat{P}(f)\right|^2 + \mu^2 f_s^2 \sum_{n=-\infty}^{+\infty} \left|\hat{P}(nf_s)\right|^2\delta(f-nf_s)}}\tag{\textasteriskcentered}
\end{equation} 